% ============================================================
% Chapter 7: Results
%   Section: RQ1 --- Finding
% Included from chapters/07-results/results.tex
% Sources: docs/experiments.md; docs/decisions/selection.md;
% backend/data/evaluation/experiments/{discovery,ranking,alignment,
% emphasis}/. Every candidate label below is the AI labelling pass's
% (Claude), not yet checked by a person.
% ============================================================

\section{RQ1 --- Finding Constructive News on Chosen Topics}
\label{sec:results-rq1}

RQ1 asks whether a retriever of the open web, scoped to chosen topics,
simplifies a journalist's task of finding constructive news on them. It
is answered by two kinds of measurement: the work the scoping saves,
and how well the ranked round orders what it finds
(Section~\ref{sec:rq-metrics}). The candidate labels used below were
assigned from title, summary and link by the AI labelling pass and have
not been checked by a person; the P1 items of
Table~\ref{tab:priorities} replace them with the author's own labels on
at least two rounds, and repeat the discovery measurement on five days.

\paragraph{The work saved.} On the first run of the discovery
experiment (5~October, two candidates per source), a round scoped to
Environment read 14 of the 34 sources then enabled ($0.41\times$ the
every-topic baseline) and 153 of 1,536 links ($0.10\times$), in 10~s
against 15~s; Culture and Health together read 23 sources
($0.68\times$) and 596 links ($0.39\times$). A person reading a scoped
round therefore starts from a tenth to two fifths of the links an
unscoped reading would offer. One run is not an interval: the planned
five days will give one.

\paragraph{The ranking separates what is worth offering.} On the 117
candidates of the five rounds recorded on 6~October, labelled before
any score existed (62 not worth offering, 33 worth offering, 22 clearly
positive), the chosen score separated stories worth offering from those
not worth it with an AUC of 0.778 [0.673, 0.868], and clearly positive
stories from the rest with 0.806 [0.679, 0.916]
(Table~\ref{tab:ranking-auc}). In the rounds, the first twenty places
of a Science and Society round held 8 stories not worth offering in
the sources' order and 1 once ranked, and 14 clearly positive instead
of 6; a Society round's first twenty went from 15 not worth offering to
6.

\begin{table}[ht]
\centering
\small
\caption{How well each signal separates the candidates of 6~October
(AUC, 95\% bootstrap interval, 86 candidates the screen keeps). Labels
by the AI labelling pass.}
\label{tab:ranking-auc}
\begin{tabularx}{\linewidth}{@{}Xcc@{}}
\toprule
\textbf{Signal} & \textbf{Worth offering vs not} & \textbf{Clearly positive vs rest} \\
\midrule
The chosen score (story 0.6, record 0.25, reliability 0.15) & 0.778 [0.673, 0.868] & 0.806 [0.679, 0.916] \\
The story alone & 0.753 [0.647, 0.850] & 0.812 [0.692, 0.917] \\
The source's record in the round & 0.684 [0.556, 0.797] & 0.664 [0.503, 0.808] \\
The source's reliability rating & 0.486 [0.356, 0.605] & 0.425 [0.294, 0.566] \\
Freshness (newer first) & 0.301 [0.195, 0.416] & 0.420 [0.269, 0.584] \\
The AI selection's score (\texttt{llama3.2:3b}) & 0.650 [0.533, 0.759] & 0.755 [0.612, 0.881] \\
\bottomrule
\end{tabularx}
\end{table}

Two signals were kept out or kept small on this evidence. Freshness
points the wrong way within the 30-day limit (AUC 0.30): the day's
breaking news is mostly the day's politics. The reliability rating
carries no signal for the mission (0.49) and has the smallest weight;
it is kept because how far a source is trusted is part of what the
publication offers. The 3-billion-parameter AI selection separated the
stories less well than the ranking on this sample.

\paragraph{Aligning the rounds with the mission and the groups.} On
9~October five one-group rounds were run live, then again after each
of four changes, and their 324 candidates labelled. Over the first
twenty places of the five rounds (100 places), pages that were not
articles went from 13 to 0, stories of another group from 21 to 13,
stories not worth offering from 29 to 17, and clearly positive stories
stayed at 34. Health, the weakest round, went from 5, 5, 10 and 3 to 0,
2, 3 and 4 on the same four counts. The format rules caught 28 of the
35 non-articles and no article; the group-fit rule told a story of the
round's group from one of another with an AUC of 0.874 [0.809, 0.930].

\paragraph{The topic emphasis.} Re-scoring every candidate of ten
labelled one-group rounds (200 places in their first twenties) at each
emphasis (Table~\ref{tab:emphasis}) shows the trade the person controls.
At the default of 0.6 the first ten places hold only the group asked
for (100 of 100) and the first twenty 194 of 200, against 184 before;
the cost is clearly positive stories of the group, 62 down to 53,
replaced by on-topic hard news that the AI selection and the person
are left to sort. The AUC for a story of the group against another's
rises from 0.706 [0.608, 0.797] to 0.825 [0.751, 0.891].

\begin{table}[ht]
\centering
\small
\caption{Ten one-group rounds re-scored at each topic emphasis: counts
over their first twenty places (200), and over their first ten (100)
in the last column. Labels by the AI labelling pass.}
\label{tab:emphasis}
\begin{tabularx}{\linewidth}{@{}lXXXXr@{}}
\toprule
\textbf{Emphasis} & \textbf{Of the group asked for} & \textbf{Worth
offering, of the group} & \textbf{Clearly positive, of the group} &
\textbf{Not worth offering} & \textbf{First ten} \\
\midrule
Before (fixed weights) & 184 & 154 & 62 & 36 & 94 \\
0 (positive impact only) & 174 & 150 & 61 & 33 & 89 \\
0.3 & 187 & 154 & 61 & 38 & 94 \\
0.5 & 192 & 150 & 53 & 45 & 98 \\
\textbf{0.6 (default)} & \textbf{194} & \textbf{150} & \textbf{53} &
\textbf{48} & \textbf{100} \\
0.7 & 194 & 139 & 46 & 60 & 100 \\
1 (topic only) & 195 & 136 & 52 & 63 & 97 \\
\bottomrule
\end{tabularx}
\end{table}

\paragraph{What this answers and what it does not.} On the evidence so
far, the scoped, screened and ranked round does what RQ1 asks of it: a
person reads a fraction of the sources and links, and the first twenty
places hold markedly fewer stories not worth offering and stories of
other groups than the sources' own order. Three limits keep this from
being a final answer. The labels are an AI's, not a person's; the
rounds come from one or two days, so the 25 rounds of 9~October are
closer to five independent samples than to 25; and order is not supply
--- a round of thirty with fourteen stories worth offering still shows
six that are not.
On the first of the five planned days, 20~October, three candidates per
source: one group alone read 27--67\% of the baseline's 33 sources and
14--29\% of its 1{,}369 links, a pair of groups 40--50\% of the links.
Environment was the lopsided group, 30 English candidates to 6 Spanish,
because Spanish feeds are filtered by title words and few Spanish titles
name the environment.
\pending{the other four days of discovery, for intervals; precision of
the first twenty and the ranking's AUC on at least two rounds labelled
by the author.}

\paragraph{A second way in: what is trending.} The trending lens
(Section~\ref{sec:design-selection}) asks RQ1 again from another side:
does starting from what is spreading find constructive stories the
topic rounds do not? A rule written before the first day fixes the
protocol (\path{backend/data/evaluation/experiments/trends/rule.md}).
On five days, a trending round over every signal is made beside the
day's topic rounds over all five groups, and the trending round's first
twenty are labelled from what a person sees, without their scores and in
a shuffled order (\texttt{python -m src.evaluation.experiments trends
run}, then \texttt{trends report}). Each is labelled worth offering or
not and clearly positive or not, with whether a constructive angle is
the story itself or only possible, and whether its subject is a private
person. Day~1, 26~October~2026, labelled by the AI pass blind to the
scores (Table~\ref{tab:trends-day1}): 12 of the first twenty were worth
offering, 5 of them clearly positive, and \textbf{none of the 12 was
offered by either topic round that day}, by link or by title similarity.
They came almost all from r/UpliftingNews (9 of its 10), against 2 of
Upworthy's 6 and 1 of Google Trends' 4. Five of the twenty had a private
person or a child as their subject, which the rules on publishing such
stories exist for.

\begin{table}[ht]
\centering
\small
\caption{The trending round of 26~October~2026, first twenty places, by
signal. Labels by the AI labelling pass, blind to the scores.}
\label{tab:trends-day1}
\begin{tabular}{@{}lrrrr@{}}
\toprule
\textbf{Signal} & \textbf{Labelled} & \textbf{Not worth} & \textbf{Worth} & \textbf{Clearly positive} \\
\midrule
r/UpliftingNews (top of the day) & 10 & 1 & 5 & 4 \\
Upworthy & 6 & 4 & 2 & 0 \\
Google Trends, Spain & 2 & 1 & 0 & 1 \\
Google Trends, United States & 2 & 2 & 0 & 0 \\
\midrule
\textbf{All} & \textbf{20} & \textbf{8} & \textbf{7} & \textbf{5} \\
\bottomrule
\end{tabular}
\end{table}

\pending{days 2--5 of the trending rounds, one a day; the precision of
the first twenty with an interval over the days; the new stories over
five days against the rule's bound of ten; whether a signal is disabled
by default; the funnel from chosen to published.}
